% ----------------------------------------------------------------------------
%
% ----------------------------------------------------------------------------

\documentclass[11pt, parskip=half]{scrartcl}       % KOMA-Skript für Artikel

%% Präambel
\usepackage[english, ngerman]{babel} % deutsche typogr. Regeln + Trenntabelle
\usepackage[T1]{fontenc}             % interner TeX-Font-Codierung
\usepackage{lmodern}                 % Font Latin Modern
\usepackage[utf8]{inputenc}          % Font-Codierung der Eingabedatei
\usepackage[babel]{csquotes}         % Anführungszeichen
\usepackage{graphicx}                % Graphiken
\usepackage{booktabs}                % Tabellen schöner
\usepackage{amsmath}      % Mathematik
\usepackage{amssymb}      % Mathematische Symbole
\usepackage{float}
\usepackage[pdftex]{hyperref}
\usepackage{subcaption}
\usepackage{url}
\hypersetup{
  bookmarksopen=true,
  bookmarksopenlevel=3,
  colorlinks,
  citecolor=blue,
  linkcolor=blue
}
\usepackage{scrhack} % unterdrückt Fehlermeldung von listings
\usepackage[backend=bibtex]{biblatex}
\addbibresource{/home/niklas/Documents/bibfile/bibliographie.bib}


%% Nummerierungstiefen
\setcounter{tocdepth}{3}     % 3 Stufen im Inhaltsverzeichnis
\setcounter{secnumdepth}{3}  % 3 Stufen in Abschnittnummerierung

\usepackage[nolist]{acronym}
\begin{acronym}
\acro{pwa}[PWA]{Progressive Web App}
\end{acronym}

\newcommand\litem[1]{\item{\bfseries#1.\space}}

\newcommand*{\N}{\mathbb{N}}
\newcommand*{\Z}{\mathbb{Z}}
\newcommand*{\Q}{\mathbb{Q}}
\newcommand*{\R}{\mathbb{R}}

% Code packages
\usepackage{listingsutf8} % Code mit UTF-8 Support
\usepackage{color,xcolor} % Farben definierbar als HTML, RGB, ...
\usepackage{textcomp}

\definecolor{alabasterGray}{HTML}{F7F7F7}
\definecolor{alabasterBlue}{HTML}{325CC0}
\definecolor{alabasterGreen}{HTML}{448C27}
\definecolor{alabasterPink}{HTML}{7A3E9D}
\definecolor{alabasterRed}{HTML}{AA3731}

\lstset{
  basewidth={0.5em,0.45em},
  extendedchars=true,
  backgroundcolor=\color{alabasterGray}, % background color (color or xcolor needed)
  basicstyle=\small\ttfamily,
  keywordstyle=\color{alabasterBlue},
  commentstyle=\color{alabasterRed},
  rulecolor=\color{black},          % frame color may change if not set (frame on comment, is comment-colored)
  stringstyle=\color{alabasterGreen},
  numberstyle=\tiny\color{gray},    %
  numbers=none,                     % where to put the line-numbers (none, left, right)
  numbersep=7pt,                    % margin between numbers and code
  stepnumber=1,                     % every n-th row will be numbered
  captionpos=b,                     % sets the caption-position to bottom
  frame=single,                     % adds a frame around the code
  keepspaces=true,                  % keeps spaces in text
  showtabs=false,                   % show tabs as special char
  showspaces=false,                 % show spaces as special char
  showstringspaces=false,           % show spaces in strings only as special char
  tabsize=2,                        % tabsize is 2 spaces
  breakatwhitespace=false,          % sets if automatic breaks should only happen at whitespace
  breaklines=true,                  % sets automatic line breaking
}
\lstset{
  literate={ö}{{\"o}}1
           {ä}{{\"a}}1
           {ü}{{\"u}}1
           % https://tex.stackexchange.com/questions/17739/listings-package-how-to-highlight-math-operators
           {true}{{{\color{alabasterPink}true}}}4
           {false}{{{\color{alabasterPink}false}}}5
           {TRUE}{{{\color{alabasterPink}TRUE}}}4
           {FALSE}{{{\color{alabasterPink}FALSE}}}5
           {True}{{{\color{alabasterPink}True}}}4
           {False}{{{\color{alabasterPink}False}}}5
}
\expandafter\def\expandafter\UrlBreaks\expandafter{\UrlBreaks\do\-\do\\/}

\begin{document}

\titlehead{\includegraphics[width=\textwidth]{src/Logo_THM_CG_FB06.png}}

%% Titelseite
\subject{Manusskript}
\title{Progressive Web Apps}
\subtitle{Kurs \enquote{Hauptseminar -- Mobile Techonologies}}
\author{Niklas Deworetzki}
\date{\today}
\maketitle

\vspace*{1.5cm}
\section*{\centerline{Zusammenfassung}}

Diese Arbeit befasst sich mit Progressive Web Applications.


\newpage

\tableofcontents
\newpage

\section{Einleitung}

https://medium.com/@amberleyjohanna/seriously-though-what-is-a-progressive-web-app-56130600a093

\acp{pwa} werden immer relevanter.
Ergänzendes Element zwischen herkömmlicher Anwendung und Website.
Was kann man damit machen? Worum geht's eigentlich?

\section{Was ist eine PWA?}

Der Begriff \acl{pwa} beschreibt eine Kategorie von Webanwendungen, welche eine Reihe an Eigenschaften erfüllen, um das Erlebnis einer nativen Anwendung in modernen Browsern zu bieten.
Geprägt wurde dieser Begriff von Alex Russell, welcher in seinem Blog\cite{russell_pwaescapingtabs} beschreibt, wie durch die Weiterentwicklung und Standardisierung von Browsern und Webtechnologien eine neue Art von Anwendung entstanden ist.
Diese Kategorie von Anwendungen entspreche dem nächsten Schritt in einer Reihe von Technologien, die versuchen, Entwicklern die Möglichkeit zu geben, mit nur einer plattformübergreifenden Anwendung auf plattformspezifische Ressourcen zuzugreifen.

Diese Kategorie von Webanwendungen ist also eine Art Nachfolger auf Technologien wie \enquote{Universal Windows Platform}-Apps\cite{msdocs_uwp} von Microsoft oder die von W3C standardisierten \enquote{Packaged Web Apps}\cite{w3c_packagedwebapps}, welche auch als \enquote{Widget} bekannt sind.
Der Begriff \enquote{Progressive} in \acl{pwa} steht für \enquote{Progressive Enhancement} (fortschreitende Verbesserung) und beschreibt ebendiese Vorgehensweise, eine Webanwendung auf möglichst vielen Geräten mit verschiedenen Kapazitäten nutzbar zu machen.
Die Idee dabei ist es, eine Grundversion der Anwendung zu erstellen, die auf jedem Gerät funktionieren kann.
Weitere Funktionalität für potentere Geräte werden dann auf diese Grundversion aufgebaut.

Da der Begriff der \ac{pwa} zusammen mit den notwendigen Technologien über einen längeren Zeitraum gewachsen ist, gibt es keine feste Definition für das, was eine \ac{pwa} ausmacht.
Die Beschreibung einer \ac{pwa} basiert vielmehr auf einer Menge an Eigenschaften, die eine Anwendung innehaben kann.
Im Folgenden werden daher einige Eigenschaften zusammengetragen und erläutert, die immer wieder für die Beschreibung von \acp{pwa} verwendet werden.
Ein besonderer Schwerpunkt wird dabei auf die Standpunkte der Google- und Mozilla-Entwickler gelegt, da diese treibende Kraft beim Entwickeln von \acp{pwa} sind und gemeinsam die Fähigkeiten dieser Gruppe von Anwendungen ausbauen.

\newpage % TODO: Eine Seite mehr an dieser Stelle. Nachher vielleicht anpassen?

Die gesammelten Eigenschaften einer \ac{pwa} lassen sich dabei in folgende Kategorien einteilen, in welcher sie auch erläutert werden.

\begin{enumerate}
  \litem{Grundlegende Eigenschaften} Eigenschaften einer \ac{pwa}, die dem Charakter einer \ac{pwa} entsprechen und ohne die ein Funktionieren der Anwendung nicht möglich wäre.

  \litem{Verbreitete Eigenschaften} Eigenschaften einer \ac{pwa}, welche von verschiedenen Quellen genannt werden, aber nicht unbedingt notwendig für das Funktionieren der Anwendung sind.

  \litem{Designmerkmale} Eigenschaften einer \ac{pwa}, die sich positiv auf das Nutzererlebnis auswirken.
\end{enumerate}


\subsection{Grundlegende Eigenschaften}

Eine der wenigen Eigenschaften einer \ac{pwa}, über die Einigkeit herrscht, ist, dass eine \ac{pwa} installierbar sein muss.
Dies bedeutet, dass es dem Nutzer möglich sein soll, die \ac{pwa} nach einem Installationsvorgang direkt vom Gerät aus starten zu können.
Ob dies nun über einen Desktopeintrag oder etwa eine Verknüpfung auf dem Startbildschirm erfolgt, ist abhängig vom Gerät des Benutzers.
Wichtig ist bloß, dass sich die installierte \ac{pwa} wie eine herkömmlich-installierte native Anwendung verhalten soll.

Dabei ist es natürlich nötig, Unterstützung vom Browser zu erhalten, welcher die \ac{pwa} startet, ausführt und anzeigt.
Dieser übernimmt auch die Verwaltung der Installation und fügt die entsprechenden Einträge in Menüs oder dem Desktop hinzu.
Um dem Nutzer ein mit einer herkömmlichen App vergleichbares Erlebnis zu bieten, erhalten \acp{pwa} beim Ausführen auch eine Sonderbehandlung vom Browser.
In Google-Chrome beispielsweise wird jede \ac{pwa} in einem eigenen Fenster gestartet, wodurch das Nutzererlebnis einer nativen Anwendung entstehen soll\cite{googledevs_pwa}.


Durch den Installationsprozess werden Teile der Anwendung lokal auf dem Gerät gespeichert.
Die Auswahl und Verwaltung der lokal gespeicherten Inhalte der Anwendung geschieht über sogenannte \enquote{Service Worker}, welche bei der Installation konfiguriert werden.
Darüber, dass das Speichern von Anwendungsinhalten durch Service Worker für eine \ac{pwa} notwendig ist, sind sich die meisten Quellen einig.
Jedoch werden verschiedene Gründe für das Speichern angeführt.
Laut dem Google-Entwickler Pete LePage optimiert das lokale Speichern von Anwendungsteilen den Startvorgang einer \ac{pwa}.
Durch das Eliminieren von Engpässen bei der Netzwerkgeschwindigkeit könne eine konsistent schnelle Ladezeit sichergestellt werden.\cite{googledev_pwaondesktop}
Bei Mozilla ist das lokale Speichern von Anwendungsinhalten unter dem Punkt der Netzwerkunabhängigkeit betrachtet.
Laut Chris David Mills solle so die Anwendung nicht nur bei schlechter Netzwerkverbindung funktionstüchtig bleiben, sondern gar ohne eine Netzwerkverbindung auskommen können\cite{mozilladevs-pwaintroduction}.

Durch die Verwendung von Service Worker ergibt sich eine weitere Eigenschaft aller \acp{pwa}.
Um Service Worker verwenden zu können, muss eine Website nämlich über % TODO


\subsection{Häufige Eigenschaften}

\subsection{Designmerkmale}


\section{Technische Voraussetzungen}

\subsection{Installationsformat}

\subsection{Datenverwaltung durch Service Worker}

\subsection{Interaktion mit Nutzern}

\subsection{Interaktion mit Hardware}

\subsection{Browserspezifische Voraussetzungen}


\section{Vorteile von PWAs}

\section{Kritik an PWAs}


\section{Fazit}

\newpage

\printbibliography

\end{document}

%%% Local Variables:
%%% mode: latex
%%% TeX-master: t
%%% End:
